\cotefile{84}
\scanpage{2}

\diagram{\begin{tikzcd}[ampersand replacement=\&]
R_1 \arrow[r,shift left,"\deleted{\unsure{\alpha}}\beta"] \& R_0 \arrow[l,shift left,"\deleted{\unsure{\beta}}\alpha"]
\end{tikzcd}}

\deleted{\unsure{hom. de monoïde|hom. de monoïdes}} \quad $R_0,R_1$ monoïdes

\[
\begin{aligned}
\alpha:R_0\longrightarrow R_1&\quad\text{hom. de monoïde}\quad (L,\delta\in\Gamma(L^{\otimes2}),T_0\in\Gamma L_0)\longmapsto (L,\Gamma)\\
\beta:R_1\longrightarrow R_0&\quad\text{\unsure{applicat. ensembl.|application ensem.|multiplicat.} (pas comp. avec produits)}\quad (L,\delta)\longmapsto(L,4\delta,0)
\end{aligned}
\]

\diagram{\begin{tikzcd}[ampersand replacement=\&]
R_\ast\arrow[d] \& {}=R_0\amalg R_1\\
\mathbb Z/2\mathbb Z \&
\end{tikzcd}}

loi de composition \inserted{$(x,y)\mapsto x\wedge y$ comm., ass., unitaire,} comp. avec $R_\ast\longrightarrow\mathbb Z/2\mathbb Z$

\[
\text{i.e.}\quad
\begin{aligned}
R_0\cdot R_0&\subset R_0\\
R_1 R_1&\subset R_0\\
R_1\cdot R_0=R_0\cdot R_1&\subset R_1
\end{aligned}
\]

\[
\begin{aligned}
x,y\in R_0&\Longrightarrow x\wedge y=x\cdot y\\
x,y\in R_1&\Longrightarrow x\wedge y=\beta(x\cdot y)\qquad\text{(pas néc. égal à $\beta(x)\,\beta(y)$)}\\
x\in R_1,\ y\in R_0&\qquad x\wedge y=y\wedge x=\deleted{\unsure{\alpha(x)y}}\,x\alpha(y)
\end{aligned}
\]

\[(x\wedge y)\wedge z=x\wedge(y\wedge z)\]

\[
\begin{aligned}
\text{a)}\quad x,y,z\in R_0&\quad\text{OK}\quad\deleted{\text{(ass. dans $R_0$)}}\qquad\text{ass. dans $R_0$}\\
\text{b)}\quad x\in R_1,\ y,z\in R_0&\quad (x.\alpha(y)).\alpha(z)=x.\underbrace{\alpha(y.z)}_{(\alpha(y)\alpha(z))}\quad\text{mult. de $\alpha$ + ass. dans $R_1$}\\
\text{c)}\quad y\in R_1,\ x,z\in R_0&\quad (\alpha(x)y)\alpha(z)=\alpha(x)(y\alpha(z))\quad\text{ass. dans $R_1$}\\
\text{d)}\quad z\in R_1,\ x,y\in R_0&\quad \alpha(x.y).z=\alpha(x)(\alpha(y).z)\quad\text{mult. de $\alpha$ + ass. de $R_1$}
\end{aligned}
\]

\[
\begin{aligned}
\text{b$'$)}\quad x\in R_0,\ y,z\in R_1&\quad\beta((\alpha(x).y).z)\overset?=x.\beta(y.z)\quad (\ast)\\
\text{c$'$)}\quad y\in R_0,\ x,z\in R_1&\quad\beta((x\alpha(y)).z)=\beta(x.(\alpha(y).z))\quad\text{ass. dans $R_1$}\\
\text{d$'$)}\quad z\in R_0,\ x,y\in R_1&\quad\underset{\Vert}{\beta(x.y)z}\overset?=\underset{\Vert}{\beta(x(y\alpha(z)))}\quad\text{comme $(\ast)$ + comm. et ass. dans $R_1$}\\
&\quad z\,\beta(y.x)\qquad\beta((\alpha(z)y)x)\\
\text{a$'$)}\quad x,y,z\in R_1&\quad\alpha(\beta(x.y)).z\overset?=x\,\alpha(\beta(y.z))
\end{aligned}
\]

Donc il faut vérifier

\[
\beta(\alpha(x).\underbrace{yz})=x\,\beta(\underbrace{yz}_{u})\qquad\text{i.e.}\tag*{\textcircled{1}}
\]

\margin{(faisant $y$ ou $z=1$)\quad $\boxed{\beta(\alpha(x)u)=x\beta(u)}$\quad $\swarrow$ $\boxed{\beta(\alpha(x))=x.\beta(1)}$}

\[
\begin{aligned}
(\alpha.\beta(xy))z&=x\,\alpha\beta(yz)\\
(\alpha\beta(x))z&=x\,(\alpha\beta(z))
\end{aligned}\tag*{\textcircled{2}}
\]

\margin{si $y=1$ cela donne et si $z=1$\quad $\boxed{\alpha\beta(x)=x.\underbrace{\alpha\beta(1)}_{\unsure{\xi|\eta}}}$}

Et \textcircled{2} devient alors $(xy\zeta)z=x(yz\zeta)$ OK par ass. et comm. de $R_1$

\diagram{\begin{tikzcd}[ampersand replacement=\&,column sep=small]
\boxed{\begin{aligned}\text{\textcircled{A}}&\quad\beta(\alpha(x)u)=x\beta(u)\\
\text{\textcircled{B}}&\quad\alpha\beta(x)=x.\underbrace{\eta}_{(\underline{\mathcal O}_X,4.\underline{1}_{\mathcal O_X})}\ \text{où }\eta=\alpha\beta(1)=\alpha(\zeta)\quad\zeta=\underbrace{\beta(1)}_{(\underline{\mathcal O}_X,4\underline1_{\mathcal O_X},0_{X_0})}\end{aligned}}\arrow[d,Rightarrow]\arrow[dr,Rightarrow] \&\\
\boxed{\beta\alpha(x)=x\zeta} \& \boxed{(\beta x)(\beta y)=\zeta\,\beta(x.y)}
\end{tikzcd}}

OK

\[
\begin{aligned}
&\deleted{\beta(1)\beta(x.y)1=\beta}\\[-2pt]
&\boxed{\underbrace{\beta(x)\beta(y)}_{\Vert\,\text{\textcircled{A}}}\overset?=\deleted{\unsure{\beta(xy)}}\underbrace{\zeta}_{\beta(1)}\cdot\beta(\overbrace{x.y}^{u})}\\
&\beta(\underbrace{\alpha\beta(x)}_{x\eta\ \text{\textcircled{B}}}y)\qquad\quad\Vert\,\text{\textcircled{A}}\\
&\beta(xy\eta)\qquad
\beta(\underbrace{\underbrace{\alpha(\zeta)}_{\eta}xy}_{\alpha\beta(xy)\ \text{\textcircled{B}}})=\beta\underbrace{\alpha\beta(xy)}_{\beta\unsure{(\beta|1\,\beta)}(xy)\eta)}
\end{aligned}
\]

Vér. \unsure{les cas|sur cas} d'espèces

\[
\begin{aligned}
\underline{x=(L,\delta)}&\qquad\beta(x)=(L,4\delta,0)\qquad\alpha\beta(x)=(L,4\delta)\\
x=(L,\delta,T_0),\ u&=(L',\delta')\qquad\qquad\quad=(L,\delta)(\underline{\mathcal O},4\underline1_X)\\
\alpha(x)&=(L,\delta)\\
\alpha(x).u&=(LL',\delta\delta')\\
\beta(\alpha(x).u)&=(LL',4\delta\delta',0)\\
\beta(u)&=(L',4\delta',0)\qquad x\beta(u)=(LL',4\delta\delta',0)
\end{aligned}
\]

$R_\ast=R\times_{\mathbb Z}\mathbb Z/2\mathbb Z$\qquad On définit

\[
\begin{aligned}
&R_\ast\xrightarrow{\deleted{\unsure{w_n}}} R_0\quad\text{par}\\
&\text{a) si }n=2\nu\qquad R_{2\nu}\xrightarrow{w_n} R_0\qquad w_n(x)=x\zeta^\nu\\
&\hspace{45pt}\deleted{\unsure{R_1}}\quad R_0\\
&\text{b) si }n=2\nu+1\qquad\deleted{n+1=2\nu+2}\qquad w_n(\unsure{x})=\beta(x)\zeta^\nu\\
&\hspace{48pt}\underset{R_1}{\underset{\Vert}{R_n}}
\end{aligned}
\]

On trouve ainsi un \uline{hom} $w:R_\ast\longrightarrow R_0$. \unsure{d'où|Un}
\[w'_1:R_\ast\longrightarrow\mathbb Z\times R_0\]

\[
\begin{aligned}
E&=\bigoplus_{1\leq i\leq n}(L_i,\delta_i)\\
w_1(E)&=\bigwedge_{1\leq i\leq n}(L_i,4\delta_i,0)
=\left(\bigotimes\unsure{\vphantom{L}\zeta} L_i,\underbrace{4^n}_{=(2^n)^2}\underbrace{\textstyle\prod\delta_i}_{\delta},0\right)
\end{aligned}
\]

\[
\begin{aligned}
\operatorname{Arf}(E)&=\left(\bigotimes L_i,\underbrace{2^\nu}_{4^{\nu'}=(2^{\nu'})^2}\underbrace{\textstyle\prod\delta_i}_{\delta},0\right)
\quad\begin{aligned}\nu&=n\ \text{si $n$ pair}\\\nu&=n+1\ \text{si $n$ impair}\end{aligned}\\
\operatorname{Arf}(E)&=w_1(E).\zeta_n\qquad \zeta_n=(\underline{\mathcal O}_X,2^n\mathbin{\unsure{:|.}}1,1
\end{aligned}
\]

\margin{de ttes façons $\nu$ est pair\quad $\nu=2\nu'$\quad $\nu'=\left[\dfrac{n+1}{2}\right]$}

\[
\begin{aligned}
w_1(\unsure{L|L_n})&=(L,4\delta,0)\\
\operatorname{Arf}(L)&=(L,4\delta,0)
\end{aligned}
\]

\margin{\deleted{On peut prendre} \inserted{$\unsure{n}$} NB si $2$ inv. \inserted{sur $X$} alors \deleted{$\unsure{d}\sim\quad\otimes$} \deleted{$\operatorname{Arf}(E)\sim$}}

\[
n-\frac{n+1}{2}=\frac{\unsure{n-1}}2\qquad
\nu''=\begin{cases}\frac n2&\text{si $n$ pair}\\ \frac{n-1}2&\text{si $n$ impair}\end{cases}
\]

\unsure{On a.}

\[
\boxed{c_1(E)=\operatorname{Arf}(E).\zeta^{(n-\nu')=\nu''}}
\qquad
\begin{aligned}
\text{où}\quad\zeta&=(\underline{\mathcal O}_X,4,0)\\
&=\operatorname{Arf}(\mu_{2X}/X)
\end{aligned}
\]

donc $\operatorname{Arf}(E)$ est un inv. plus fin que $c_1(E)$ (si $2$ \uline{non} \uline{inv.} sur $X$)

\[
\begin{aligned}
c_1(E)&\simeq\left(\bigotimes L_i,\unsure{\delta},0\right)\\
\operatorname{Arf}(E)&\simeq\left(\bigotimes L_i,\delta,0\right)
\end{aligned}
\]

Mais $c_1(E)$
% archive sheet number: 1
